\label{chap:discussion}

\section{Vincoli scacchistici e spazio decisionale}

La progressione da A ad A3 mostra che la rappresentazione della decisione è
importante almeno quanto l'encoder condiviso. Il Modello A può assegnare probabilità
a classi incompatibili con la posizione corrente; A1 e A2 riducono questo problema
filtrando o mascherando le mosse illegali, mentre A3 costruisce direttamente il
confronto tra le sole mosse legali disponibili. Sul test Lichess, il passaggio dal
39,62\% di A al 49,81\% di A1 isola il beneficio del vincolo in inferenza senza
introdurre un nuovo checkpoint. A3 raggiunge il 67,56\%, suggerendo che formulare
l'apprendimento come ranking contestuale delle candidate sia più efficace del solo
filtro applicato a un vocabolario globale.

Questa evidenza non implica che la legalità esaurisca la difficoltà scacchistica.
Il generatore di mosse elimina azioni impossibili, ma il modello deve ancora
distinguere la continuazione tatticamente corretta dalle alternative legali. Il
risultato indica piuttosto che incorporare un vincolo certo del dominio riduce uno
spazio decisionale inutilmente ampio e allinea meglio l'obiettivo di addestramento
con quello di valutazione.

\section{Reranking post-mossa e limite di recupero}

A4 aggiunge una seconda prospettiva: non valuta soltanto la posizione corrente e
la coppia origine--destinazione, ma costruisce il grafo ottenuto dopo ogni candidata
e ne usa il contesto per il reranking. L'aumento della Top-1 dal 67,56\% di A3
all'85,41\% su Lichess è coerente con l'utilità di osservare le conseguenze
immediate della mossa. Sul benchmark classico il vantaggio descrittivo è più
contenuto, dall'11,0\% al 15,5\%, e il test appaiato non raggiunge la soglia
convenzionale del 5\% (\(p\approx0{,}0931\)); il campione di 200 problemi limita la
precisione della stima.

Il risultato di A4 è condizionato dallo stadio di recupero. Il reranker riceve solo
le cinque migliori candidate di A3: se la mossa accettata non compare in questo
insieme, A4 non può selezionarla. La Top-5 di A3 costituisce quindi un tetto di
recupero end-to-end; su YACPDB essa è pari al 34,5\%. A4 migliora l'ordinamento
all'interno dell'insieme disponibile, ma non risolve gli errori di recupero a monte.
Questa distinzione evita di attribuire al solo encoder post-mossa l'intera
prestazione della pipeline a due stadi.

\section{Timing e confronto con i modelli linguistici}

Nel confronto controllato il Modello B perde 1,58 punti percentuali di Top-1
rispetto ad A3; i casi discordanti e il test esatto di McNemar indicano che la
differenza non è spiegata soltanto da oscillazioni casuali nel campione valutato.
Le ablazioni neutralizzate restano inferiori ad A3. L'evidenza supporta quindi una
conclusione circoscritta: la codifica temporale sintetica implementata non migliora
lo scorer di mosse legali.

Non segue che il tempo di riflessione sia irrilevante in generale. I valori usati
sono generati deterministicamente e condizionati dal rating; non misurano il tempo
impiegato da una persona per risolvere quello specifico puzzle. È plausibile, ma non
dimostrato dagli esperimenti, che il segnale non aggiunga informazione indipendente
o che l'encoder non ne sfrutti la forma. Verificare queste ipotesi richiederebbe
tempi reali allineati ai puzzle o un diverso disegno sperimentale.

Nel protocollo rigoroso UCI, entrambi i modelli Qwen ottengono accuratezze molto
inferiori ai modelli a grafo a ogni profondità Lichess; RQ1 riceve dunque una
risposta uniforme nei cinque gruppi disponibili. Il risultato dimostra l'efficacia
dei sistemi completi sotto il protocollo congelato, non una superiorita generale
delle GNN su ogni uso possibile dei modelli linguistici. Prompt singolo, generazione
di una sola mossa, assenza di strumenti scacchistici e parser rigoroso delimitano il
confronto. L'elevata quota di output non interpretabili o illegali è parte della
prestazione end-to-end misurata, ma non equivale da sola a una misura pura del
ragionamento scacchistico interno.

\section{Generalizzazione e limiti sperimentali}

Il calo di A3 e A4 da Lichess a YACPDB segnala un cambio di distribuzione: le
composizioni classiche differiscono dai puzzle estratti da partite per provenienza,
costruzione e profondità. I dati non provano che YACPDB sia intrinsecamente più
difficile; mostrano invece che le prestazioni sulla distribuzione principale non si
trasferiscono integralmente. Il campione di 200 problemi, 20 per profondità, rende
inoltre instabili i confronti fini tra categorie.

La portata del lavoro è circoscritta da ulteriori limiti. GraphMate predice la
prima mossa risolutiva, non l'intera variante. Il timing sintetico manca in YACPDB,
quindi RQ2 resta limitata a Lichess; A4 dipende dal recupero Top-5 di A3. Le
baseline linguistiche coprono un protocollo locale, non ogni modello, prompt o
integrazione con motori. Checkpoint e budget sono quelli delle configurazioni
ufficiali: i risultati descrivono esecuzioni congelate, non limiti teorici.
